% L15a strategy deck: the correlation-gamma picking rule.
% One short deck per rule on the strategy menu, read next to the rule's code file
% (strategies/picking/correlation_gamma.jl). Written by make_strategy_decks.py (session scratchpad, 2026-11-27) and then
% kept by hand: edit this file, not the generator. Build with make in this folder: LaTeX runs from strategies/,
% so paths are relative to it.
\documentclass[aspectratio=169,11pt]{beamer}
\usepackage{vnslides}

\newcommand{\repo}{https://github.com/varnerlab/CHEME-5660-CourseRepository-Fall-2026/blob/main/lectures/week-15/L15a}

\title{\texorpdfstring{{\vnheavy L15a}}{L15a}: Correlation Gamma}
\subtitle{Picking Rule, CHEME 5660 Strategy Menu, Cornell University}
\author{Copyright Jeffrey Varner 2026. All Rights Reserved}

\begin{document}

\vntitlepage

\begin{frame}{Disclaimer and Risks}
  \vspace{0.6em}
  \textbf{\ink{Training only.}} This content is for training and informational
  purposes only. The teaching team does not make, give, or endorse any offer or
  solicitation to buy or sell securities or derivative products, or provide
  investment or trading advice or strategy.

  \vspace{1.1em}
  \textbf{\ink{Risk.}} Trading involves risk and is generally not recommended for
  those with limited resources, experience, or a low risk tolerance. Past
  performance does not guarantee future success.

  \vspace{1.1em}
  \textbf{\ink{Responsibility.}} You are responsible for your investment decisions
  based on your financial circumstances, objectives, risk tolerance, and
  liquidity needs.
\end{frame}

\begin{frame}{The Idea}
  \vspace{0.3em}
  Tickers with high preference weights can move together, such as stocks in the same industry. The
  correlation-gamma rule goes down the ranking of preference weights and skips a ticker that moves too closely
  with a ticker already taken:
  \begin{itemize}\setlength{\itemsep}{0.5em}
    \item \hilite{Universe}: 491 stocks and 14 exchange-traded funds (ETFs), without \texttt{SPY} and the Treasury bill fund \texttt{SGOV}.
    \item \hilite{Limit}: the correlation of daily growth rates with each ticker already taken is at most 0.7, over the last 252 trading days.
    \item \hilite{Size}: the basket holds $K=13$ tickers at every decision.
  \end{itemize}

  \vspace{0.5em}
  The engine runs the rule at each decision: on the start day and whenever a trigger fires. Between decisions, the basket stays the same.
\end{frame}

\begin{frame}{The Preference Weight}
  \vspace{0.3em}
  At each close $t$, the engine scores every ticker $i$ with its preference weight $\gamma_{i}(t)$:
  {\setlength{\jot}{3pt}
  \begin{align*}
    \tilde{g}_{i,t} &= \alpha_{i}(t)+\beta_{i}(t)\,\tilde{g}_{M,t} && \text{the expected growth rate (1/yr)}\\
    \gamma_{i}(t) &= \tanh\left(h\,\frac{\tilde{g}_{i,t}}{\max\left(|\beta_{i}(t)|,\,\beta_{\text{floor}}\right)^{\xi_{t}}}\right)
      && h=1\text{ year},\ \beta_{\text{floor}}=0.1
  \end{align*}}%
  where the single index model gives $\alpha_{i}(t)$ and $\beta_{i}(t)$, $\tilde{g}_{M,t}$ is the 252-day
  exponential moving average of market growth, and the market state $\xi_{t}$ is positive when the market falls.

  \vspace{0.5em}
  The preference weight lies between $-1$ and $1$, with the sign of the expected growth rate.
\end{frame}

\begin{frame}{How the Market State Tilts the Preference Weight}
  \vspace{0.3em}
  The market state $\xi_{t}$ sets how the market sensitivity scales the expected growth rate:
  \begin{center}
    \includegraphics{slides/figs/Fig-L15a-Deck-Tilt.pdf}
  \end{center}
  At the same positive expected growth rate, the preference weight favors a low market sensitivity when the
  market falls and a high one when the market rises.
\end{frame}

\begin{frame}{The Rule}
  \vspace{0.3em}
  On decision day $t$, let $\rho_{ij}(t)$ be the correlation of the daily growth rates of tickers $i$ and $j$
  over the 252 trading days through day $t$. Start with an empty basket $\mathcal{B}(t)$:
  \begin{enumerate}\setlength{\itemsep}{0.5em}
    \item Rank every ticker by $\gamma_{i}(t)$, from highest to lowest.
    \item Go down the ranking. Take ticker $i$ when $\rho_{ij}(t)\leq 0.7$ for every ticker $j$ already in
      $\mathcal{B}(t)$, and stop when $\mathcal{B}(t)$ holds $K=13$ tickers.
    \item If fewer than 13 tickers pass, fill the basket with the highest-ranked tickers not yet taken.
  \end{enumerate}
\end{frame}

\begin{frame}{The Basket of a Made-Up Ranking}
  \vspace{0.3em}
  The rule goes down a made-up ranking and skips each ticker with $\rho_{ij}(t)>0.7$ for a ticker $j$ already in
  the basket:
  \begin{center}
    \includegraphics{slides/figs/Fig-L15a-Deck-Picks-Correlation.pdf}
  \end{center}
  Each arc joins a skipped ticker to a ticker already taken, and its label gives their correlation. After four
  skips, the 13th ticker taken is at rank 17.
\end{frame}

\begin{frame}{Reading the Code}
  \vspace{0.3em}
  The rule's code file is
  \href{\repo/strategies/picking/correlation_gamma.jl}{\texttt{strategies/picking/correlation\_gamma.jl}}:
  \begin{itemize}\setlength{\itemsep}{0.5em}
    \item \texttt{pick\_correlation\_gamma(tickers, gamma, growth; count, max\_correlation)} goes down the ranking and fills a short basket. Each column of \texttt{growth} holds one ticker's daily growth rates.
    \item \texttt{MyL15aCorrelationGammaScreen} is the type the engine calls. It holds the settings \texttt{count~=~13}, \texttt{max\_correlation~=~0.7}, and \texttt{window~=~252}.
    \item \texttt{MPC.pick(screen, md, t; $\gamma$\_t, ...)} computes the growth rates over the window, runs the rule with the day's preference weights \texttt{$\gamma$\_t}, and returns the basket's columns in the market data.
  \end{itemize}

  \vspace{0.5em}
  The \href{\repo/strategies/CHEME-5660-L15a-Strategy-Menu-Fall-2026.ipynb}{strategy menu} lists every rule, and the \href{\repo/engine/CHEME-5660-L15a-Trading-Engine-Fall-2026.ipynb}{trading engine} describes the preference weight and the triggers.
\end{frame}

\end{document}
